% Nature style: the main text opens without a heading. npj: "Avoid
% 'Introduction' as a heading." Do not add one.


A chest-radiograph classifier that knows when to defer is more useful in a
hospital than one that is merely accurate. Selective prediction formalises
this: the model returns a decision only when its confidence clears a threshold,
and refers the remainder to a human reader. The threshold fixes an operating
point --- a coverage --- and with it the workload the system imposes on
radiology.

That operating point is chosen once, on data from the hospital where the system
was developed. It is then deployed somewhere else. Whether it survives the move
is an empirical question that is rarely asked, and the way it is usually asked
answers something narrower than intended: models are re-evaluated at a second
site, but the abstention threshold is re-fitted there too. A re-fitted threshold
tests whether a \emph{method} transfers. It cannot test whether a \emph{deployed
policy} transfers, because the policy that would actually ship is the frozen
one.

A second shift travels with the first and is usually left implicit. Multimodal
chest-radiograph models increasingly consume clinical context alongside the
image --- the indication, the referring history, the reason the study was
ordered. This text is not a stable input. Its availability is a property of how
a particular hospital documents care, and it varies across institutions in ways
that have nothing to do with the patient. A model that learned to lean on the
indication field at one site may find it thin, absent, or differently written at
the next.

Institution and context therefore shift together. Studies that vary one hold the
other fixed: cross-institution work is typically image-only, and multimodal
context work is typically single-institution. Holding one fixed cannot reveal
whether they interact, and interaction is precisely what determines deployment
risk. If a multimodal model is more accurate at the source site \emph{because}
of context, and context is scarcer at the target site, then the measured benefit
and the unmeasured fragility are the same phenomenon observed from two sides.

This paper contributes an evaluation protocol that crosses the two shifts
deliberately, so their interaction can be estimated rather than assumed away:

\begin{itemize}
  \item \textbf{Two sites.} A source hospital, where models are developed and
    the selective policy is calibrated, and an external hospital used only for
    evaluation.
  \item \textbf{Three controlled context states at each site.} Context left
    intact, context removed, and context replaced with another patient's ---
    the last separating \emph{absence} of information from \emph{wrong}
    information.
  \item \textbf{A frozen policy.} Per-pathology thresholds and abstention
    cutoffs are selected once at the source site and transferred unchanged, so
    what is evaluated is the policy that would actually ship.
  \item \textbf{Preregistration} in a machine-readable, hash-manifested
    protocol, with every amendment dated and its confirmatory status recorded,
    so the boundary between prespecified and post-hoc is auditable rather than
    asserted.
\end{itemize}

The contribution is the protocol and what it measures. We propose no new
architecture and no new uncertainty algorithm; the backbones are deliberately
conventional, held constant so that architecture cannot serve as a competing
explanation for any effect observed. This is not a prospective clinical study or
a regulatory validation.

Two further contributions emerged from building the protocol and are reported
because they generalise beyond it. First, a structural hazard in deriving
pre-diagnostic context from radiology reports: a substantial minority of reports
embed a preliminary interpretation inside the report body, so any pipeline that
takes ``everything before the findings section'' as context silently ingests
post-diagnostic text (\cref{sec:report-structure}). Second, a property of the
decision rule this literature commonly uses --- requiring both that a confidence
interval exclude zero and that the point estimate reach a materiality threshold
--- which caps statistical power at \SI{50}{\percent} when the true effect
equals that threshold, regardless of sample size (\cref{sec:power}).

\label{sec:related-work}

Every work cited here was verified against its official publication page or
full text.

\paragraph{Selective prediction under distribution shift}

Selective classification has been extended to distribution shift outside
medicine. The closest methodological work generalises selective classification
so that a model rejects not only ambiguous in-distribution inputs but
label-shifted and covariate-shifted ones, using margin-based confidence scores
requiring no access to training data, evaluated on ImageNet, ImageNet-C,
OpenImage-O, iWildCam, Amazon and CIFAR~\cite{selective_shift_2025}.

That line establishes the machinery this paper applies. It does not take it to a
clinical setting, does not involve a second hospital, and does not transfer a
threshold frozen at one site to another. Whether an abstention policy remains at
its intended operating point after institutional transfer is not addressed
there, and is the question we ask.

\paragraph{Clinical context as a model input}
\label{sec:rw-context}

Clinical history improves chest-radiograph classification. The most directly
comparable work fuses view-specific image encoders with clinical history through
an attention-refinement mechanism on MIMIC-CXR-JPG, reporting
\num{135682} images with clinical history alongside a further \num{59846}
without~\cite{yang2025mcxnet}.

That paper is a premise for ours rather than a competitor, and the distinction
is worth stating precisely. It is single-institution, evaluated on held-out
splits of one source, with no abstention policy, no audit of whether the
diagnosis leaks into the history text, and no calibration transfer. Its
\num{59846} context-free images are an \emph{observation} of naturally occurring
missingness. We \emph{intervene} on context under a frozen informativeness rule,
which is what allows an effect to be attributed to the manipulation rather than
to whatever else distinguishes patients whose indication field happens to be
empty.

A second study reaches the same conclusion by a different route, fusing
multi-view radiographs with clinical indications and vital signs and reporting
an ablation in which each added modality
contributes~\cite{sloan2025clinically}. That ablation is additive: it measures
what each modality contributes as the model is built. Ours is interventional: it
measures what a fixed, already-trained model loses when context degrades at
deployment. The distinction is not pedantic, since a component that improves a
model during development need not be one whose absence degrades it
proportionally afterwards, and only the second quantity bears on deployment
risk. That work also evaluates on benchmark splits of a single source, with no
second institution, no abstention, and no threshold transfer.

\paragraph{Missing modalities as an architectural problem}

A body of work treats missing modalities as something to engineer around.
Transformer-based bi-modal fusion modules combined into a tri-modal framework,
trained with multivariate losses for robustness to missing modalities on a
14-label diagnosis task over MIMIC-IV and MIMIC-CXR, is
representative~\cite{wang2025missing}; so is a mixture-of-experts framework
selecting experts according to which modalities are present, for mortality,
length-of-stay and readmission prediction~\cite{wang2025moehealth}.

The orientation differs from ours in a way that matters. There, missingness is
an obstacle and robustness is the objective; the natural question is how to keep
performance up when a modality drops out. Here, missingness is the
\emph{estimand}. We do not ask how to build a model that resists context loss,
but how much reliability a conventionally built model loses when context
degrades, and whether that loss is larger at a hospital the model has never
seen. Both approaches also evaluate within a single source, so neither can speak
to transfer.

\paragraph{Cross-institution generalisation in chest radiography}

Cross-domain generalisation for chest-radiograph models is well studied.
Foundational work across multiple datasets argues that cross-domain failure is
driven by label shift rather than image shift, observing that models with strong
individual performance disagree substantially with one
another~\cite{cohen2020limits}. More recent multi-center benchmarking over
\num{145000} images from PadChest and NIH Chest X-ray evaluates head-versus-tail
performance, calibration and cross-center generalisation gaps, and concludes
that detecting rare findings under multi-center shift remains
challenging~\cite{cxrlt2026}.

A related strand examines robustness to technical acquisition parameters rather
than to institution. Comparing four architectures for the sensitivity gap
between anteroposterior and posteroanterior views, model rankings reverse
completely on a second dataset: a chest-radiograph-specific vision-language
model with the smallest internal gap (\SI{14.3}{\percent}) shows the largest
external one (\SI{48.9}{\percent}), while a self-supervised model stays near
its internal value~\cite{farquhar2026acquisition}. That the second dataset
shares an institution with the first, differing principally in labelling
convention, makes the reversal more striking rather than less: internal ranking
failed to predict external ranking without a change of hospital at all.

This line is image-only. It varies institution, acquisition, or labelling while
holding input availability fixed --- the mirror image of
\cref{sec:rw-context} --- and none of it evaluates abstention or transfers a
threshold frozen. Thresholds are re-fitted per dataset, typically by Youden's
statistic, so what is measured is whether a method can be re-tuned rather than
whether a policy survives the move. That the most recent multi-center benchmark
still names calibration and cross-center generalisation as open problems
supports treating the crossed question as unresolved rather than answered
piecemeal.

\paragraph{The gap}

Each axis has been studied in isolation: selective prediction under shift
outside medicine~\cite{selective_shift_2025}; clinical context helping within a
single hospital~\cite{yang2025mcxnet,sloan2025clinically}; missing modalities
handled architecturally within a single
source~\cite{wang2025missing,wang2025moehealth}; institutional, acquisition and
labelling shift characterised for image-only
models~\cite{cohen2020limits,cxrlt2026,farquhar2026acquisition}.

Absent is the crossing: an evaluation in which institution and pre-diagnostic
context availability vary together, under a selective policy frozen at the
source site and transferred without refitting, with the interaction estimated
rather than assumed. That is the design this paper specifies and reports.
